% Generated draft: requires mathematical and source review.
\documentclass[11pt]{article}
\usepackage[margin=25mm]{geometry}
\usepackage{fontspec}
\setmainfont{Latin Modern Roman}
\usepackage{amsmath,amssymb,mathtools,unicode-math}
\setmathfont{Latin Modern Math}
\usepackage{enumitem,needspace,xurl,hyperref}
\setlength{\parindent}{0pt}
\setlength{\parskip}{5pt}
\setlength{\emergencystretch}{3em}
\newcommand{\N}{\mathbb{N}}
\newcommand{\Z}{\mathbb{Z}}
\newcommand{\Q}{\mathbb{Q}}
\newcommand{\R}{\mathbb{R}}
\newcommand{\C}{\mathbb{C}}
\newcommand{\F}{\mathbb{F}}
\newcommand{\A}{\mathbb{A}}
\newcommand{\PP}{\mathbb{P}}
\newcommand{\E}{\mathbb{E}}
\newcommand{\eps}{\varepsilon}
\newcommand{\dd}{\,\mathrm{d}}
\begin{document}
\section*{UnsolvedMath ID: 20000442}
\textbf{Source problem code:} catalogue entry
\subsection{Definitions and mathematical statement}
For a proper variety $X/k$, direct $R$-equivalence links two $k$-points by a rational curve parametrized over $k$; $R$-equivalence is generated by chains of such links. Write $X(k)/R$ for the resulting set. In the intended smooth proper geometrically integral characteristic-zero setting, rational connectedness is geometric; over $\C$ the source describes it by $X(\C)/R$ being a singleton.

\textbf{Original question (Colliot-Thélène).} What number-field results follow from techniques for rationally connected varieties? As background the source attributes to Kollár finiteness of $X(k)/R$ over local fields for rationally connected $X$, and to Kollár--Szabó triviality of $X(k_v)/R$ for almost all places of a number field. Exact smoothness/properness hypotheses and bibliography are not supplied in the short account.

Ellenberg's original remark asks whether, for a rationally connected variety over a number field, some extension $K\supset k$ makes $X(L)/R$ a singleton for every $L\supset K$. Raskind reports a negative answer using a conic bundle over $\PP^2$, without specifying it. The quantifiers do not say whether $K$ or $L$ must be finite/algebraic extensions; universal triviality over arbitrary field extensions is stronger than a question restricted to number-field extensions.

\textbf{Attached unreviewed Brauer obstruction.} For every smooth proper geometrically integral $X$ over a characteristic-zero field with $\operatorname{Br}(X_{\overline k})\ne0$, the aid asserts: for every field extension $K/k$ there is a finite extension $E/K$, a point $P\in X(E)$, and the field $L=E(X_E)$ such that the generic point $\eta\in X(L)$ and the constant point $P_L$ are not $R$-equivalent. Here $\operatorname{Br}$ is the cohomological Brauer group, and the nonzero hypothesis is geometric, not merely a constant class from $\operatorname{Br}(k)$. The proposed witness arises from a descended geometric Brauer class and its evaluations. This prevents universal $R$-triviality after every proposed base extension in the all-field-extension sense. The witness $L$ is a function field and is not claimed finite over $K$, so it does not alone settle the finite-number-field-only interpretation.

Separately, for a smooth proper rationally connected variety over a number field, the aid combines the reported local results to obtain finiteness of
\[\prod_vX(k_v)/R.\]
If $X(k)/R$ were infinite, its localization map to this finite product would have an infinite fiber. Thus infinitely many global classes would be indistinguishable by their local $R$-classes at every completion. This is a conditional consequence, not a construction of an infinite global quotient or a proof of a local-to-global principle. All propositions and status claims are recorded as supplied research-aid assertions requiring review.
\subsection{Short English statement}
The source asks for arithmetic consequences of rational connectedness and an eventual universal $R$-triviality question. The aid gives a geometric Brauer obstruction using a function-field generic point, and observes that an infinite global $R$-quotient would have an infinite fiber in the finite product of local quotients. Arbitrary field extensions and finite number-field extensions remain distinct.
\subsection{Sources}
\begin{itemize}
\item Local UnsolvedMath registry, permanent ID 20000442, source code catalogue entry; \texttt{aim-algebraic-number-theory-notes.json}, notes[93], tag \texttt{question}.
\item AIM workshop \emph{Rational and integral points on higher dimensional varieties}, Question 10 (Colliot-Thélène): \url{https://aimath.org/WWN/qptsurface2/qptsurface2.pdf}. Original background names Kollár and Kollár--Szabó; the remark names Ellenberg and Raskind. No publication locators or conic-bundle example are supplied.
\item Supplied verification explains OCR correction of the complex embedding, Szabó's name and $\PP^2$; it refers to HTML/PDF comparison but this draft claims no source visit.
\item Attached machine-generated research attempt 1: result \texttt{partial\_result}, problem status at run \texttt{open}, rigor label \texttt{proved}, novelty confidence low. It is not peer reviewed and requires independent verification before citation. Original literature field empty.
\end{itemize}
\end{document}
